% ====================================================================
% Chapter 4: Transfer Learning
% (Digitized directly from Author Manuscript Pages 106 - 117)
% ====================================================================

\chapter{Transfer Learning}

A.I. systems learn by extracting patterns from data. In traditional machine learning and deep learning, every new task is usually learned from scratch. This means that whenever a new problem arises, the model begins with randomly initialized parameters and must learn everything again, even if it has previously solved similar problems.

\begin{propertybox}[Definition: Transfer Learning]
\textbf{Transfer Learning} is the process of improving learning in a target task by reusing knowledge acquired from a previously learned source task.
\end{propertybox}

\section{Why Is Transfer Learning Needed?}

Early deep learning systems faced three major challenges:

\subsection{Data Scarcity}
Deep neural networks contain millions or even billions of trainable parameters. For example:
\begin{itemize}
    \item \textbf{LeNet:} $60,000$ parameters
    \item \textbf{AlexNet:} $60$ million parameters
    \item \textbf{ResNet:} $25$ million parameters
    \item \textbf{GPT-3:} $175$ billion parameters
\end{itemize}
Training such models requires enormous datasets.

\subsection{Computational Cost}
Training a deep neural network from random initialization is computationally expensive. Training large models may require:
\begin{itemize}
    \item Multiple GPUs / TPUs
    \item Days or weeks of computation
    \item Large electrical power consumption
\end{itemize}
Transfer Learning significantly reduces these costs because much of the learning has already been completed during pretraining.

\subsection{Training Time}
Training time often decreases from weeks $\to$ days $\to$ hours $\to$ sometimes only minutes!

\begin{propertybox}[Human Analogy of Transfer Learning]
Transfer Learning behaves like human learning and can be viewed as:
\begin{equation}
\text{Knowledge}_{\text{new}} = \text{Knowledge}_{\text{previous}} + \text{New Experience}
\end{equation}
Transfer Learning attempts to reproduce this behavior in artificial neural networks.
\end{propertybox}

\section{Types of Transfer Learning}

Transfer Learning is classified according to whether the source and target tasks/domains are the same or different.

\subsection{Inductive Transfer Learning}
The target task is different from the source task, regardless of whether the domains are the same or different ($T_S \neq T_T$).
\begin{itemize}
    \item \textbf{Source:} ImageNet $\to$ Image Classification
    \item \textbf{Target:} Medical Image Classification
\end{itemize}
Although both involve images, the classification tasks are different.

\subsection{Transductive Transfer Learning (Domain Adaptation)}
The task remains the same, but the data distribution (domain) changes ($D_S \neq D_T, T_S = T_T$).
\begin{itemize}
    \item \textbf{Trains on:} Daytime street images
    \item \textbf{Test on:} Nighttime street images
\end{itemize}
The task (object detection) is identical, but the image distributions differ.

\subsection{Unsupervised Transfer Learning}
Both the source and target tasks are unsupervised, meaning there are no labels. Typical tasks include:
\begin{itemize}
    \item Clustering
    \item Dimensionality Reduction
    \item Representation Learning
\end{itemize}
\textit{Example:} A self-supervised model trained on millions of unlabeled images is transferred to another unlabeled image dataset.

\section{Transfer Dynamics: Positive vs. Negative Transfer}

\begin{advbox}[Positive Transfer]
Knowledge learned from the source task improves performance on the target task:
\begin{equation}
\text{Accuracy}_{\text{Transfer}} > \text{Accuracy}_{\text{scratch}}
\end{equation}
\textit{Example:} ImageNet $\to$ Medical X-ray. The pretrained CNN already knows edges, textures, and shapes, allowing faster learning and higher accuracy.
\end{advbox}

\begin{advbox}[Negative Transfer]
Transferred knowledge actually reduces performance:
\begin{equation}
\text{Accuracy}_{\text{Transfer}} < \text{Accuracy}_{\text{scratch}}
\end{equation}
\textit{Example:} A CNN trained on satellite images is transferred to handwritten digit recognition. The learned features are largely unrelated, so the transferred knowledge may hurt performance.
\end{advbox}

\subsection{When Transfer Learning Works Well}
Transfer Learning works well when:
\begin{itemize}
    \item Source and target domains are similar.
    \item The pretrained model learned general features.
    \item The target dataset is small.
    \item The source dataset is large and diverse.
    \item The target task is related to the source task.
\end{itemize}

\textbf{Common Examples:}
\begin{itemize}
    \item ImageNet $\to$ Medical Images
    \item ImageNet $\to$ Plant Disease Detection
    \item BERT $\to$ Sentiment Analysis
    \item GPT $\to$ Question Answering
\end{itemize}

\subsection{When Transfer Learning May Fail}
Transfer Learning may fail when:
\begin{itemize}
    \item Source and target domains are very different.
    \item The pretrained features are not useful.
    \item The source model contains strong biases.
    \item The target dataset is extremely different.
    \item Too much fine-tuning causes catastrophic forgetting.
\end{itemize}

\section{Application Strategies: Feature Extraction vs. Fine-Tuning}

Transfer Learning can be applied in two main ways:

\subsection{Feature Extraction}
Uses the pretrained model as a fixed feature extractor. The pretrained layers are frozen, meaning their weights remain unchanged during training. Only a new output layer (classifier) is added and trained on the target dataset (we just replace the last layer of the network).

\textbf{Characteristics:}
\begin{itemize}
    \item Pretrained weights are frozen.
    \item Only the new classifier learns.
    \item Very fast training.
    \item Requires a small dataset.
    \item Low risk of overfitting.
\end{itemize}

\subsection{Fine-Tuning}
Fine-Tuning starts with a pretrained model but allows some or all of the pretrained weights to continue learning. Instead of freezing every layer, selected layers are unfrozen and updated using the new dataset.

\textbf{Characteristics:}
\begin{itemize}
    \item Some or all pretrained layers are trainable.
    \item The classifier is always replaced and trained.
    \item Produces better adaptation to the target task.
    \item Requires more data.
    \item Higher computational cost.
    \item Greater risk of overfitting if the dataset is too small.
\end{itemize}

\begin{propertybox}[Do We Change the Input Layer in Transfer Learning?]
\textbf{Usually no.} The input layer is kept because the pretrained model expects a specific input size:
\begin{itemize}
    \item \textbf{VGG16:} $224 \times 224 \times 3$
    \item \textbf{ResNet50:} $224 \times 224 \times 3$
    \item \textbf{EfficientNetB0:} $224 \times 224 \times 3$
\end{itemize}
If our images have a different size, we usually resize the images, not the network.
\end{propertybox}

\section{Advanced Fine-Tuning Techniques}

\subsection{Partial Fine-Tuning}
Instead of training the entire network, only the last few layers are unfrozen while the earlier layers remain frozen. This is the most commonly used fine-tuning strategy because early layers learn generic features while deeper layers learn task-specific representations.

\subsection{Progressive Unfreezing}
Instead of unfreezing many layers at once, layers are unfrozen gradually:
\begin{enumerate}
    \item Classifier (last layer)
    \item Classifier + Last block
    \item Classifier + Last two blocks
    \item Classifier + Entire Network
\end{enumerate}
This provides more stable training and reduces catastrophic forgetting.

\subsection{Differential Learning Rates}
Different layers receive different learning rates. Early layers receive very small updates because they already contain useful general features, while newly added layers receive larger updates.

\begin{table}[htbp]
\centering
\begin{tabular}{l | l}
\hline
\textbf{Layer} & \textbf{Learning Rate} \\
\hline
Early layers & $1 \times 10^{-5}$ \\
Middle layers & $1 \times 10^{-5}$ \\
Last layers & $1 \times 10^{-4}$ \\
New Classifier & $1 \times 10^{-3}$ \\
\hline
\end{tabular}
\caption{Differential Learning Rates assigned per network layer depth.}
\end{table}

\section{Catastrophic Forgetting \& Regularization}

\textbf{Catastrophic Forgetting} occurs when fine-tuning causes the pretrained model to lose previously learned useful knowledge.

\begin{itemize}
    \item \textbf{Common Causes:} High learning rate, very small dataset, or training all layers immediately.
    \item \textbf{Solutions:} Freeze early layers, use small learning rates, apply progressive unfreezing, and use regularization.
\end{itemize}

\subsection{Regularization Techniques}
Regularization includes techniques that reduce overfitting by preventing the model from becoming too specialized on the training data. Common methods include Dropout, Weight Decay, Early Stopping, Data Augmentation, and Label Smoothing. The objective is to improve generalization on unseen data.

\begin{itemize}
    \item \textbf{Weight Decay ($L_2$ Regularization):} Penalizes large weights by adding a regularization term to the loss function:
    \begin{equation}
    L_{\text{total}} = L_{\text{task}} + \lambda \sum_i w_i^2
    \end{equation}
    
    \item \textbf{Early Stopping:} Instead of training for a fixed number of epochs, training ends once improvement has stalled for a predefined number of epochs (called patience).
    \begin{itemize}
        \item Prevents overfitting.
        \item Saves training time.
        \item Keeps the best-performing model.
    \end{itemize}
\end{itemize}

\subsection{Batch Normalization Issues}
During fine-tuning on small datasets, updating Batch Normalization statistics (mean and variance) can make training unstable. A common practice is:
\begin{itemize}
    \item Freeze BatchNorm layers when the target dataset is small.
    \item Update them only if sufficient data is available.
\end{itemize}

\section{Practical Recipes}

\begin{propertybox}[Scenario 1: Very Small Dataset (< 1000 images)]
\begin{itemize}
    \item Freeze the backbone.
    \item Train only the new classifier.
    \item Small learning rate.
    \item Use data augmentation.
\end{itemize}
\end{propertybox}

\begin{propertybox}[Scenario 2: Medium Dataset (1000 to 10000 images)]
\begin{itemize}
    \item Train the new classifier.
    \item Fine-tune the last few blocks.
    \item Use differential learning rates.
    \item Apply weight decay and early stopping.
\end{itemize}
\end{propertybox}

\begin{propertybox}[Scenario 3: Large Dataset (> 10000 to 50000 images)]
\begin{itemize}
    \item Fine-tune most or all layers.
    \item Use progressive unfreezing if needed.
    \item Tune learning rates carefully.
    \item Monitor validation performance to avoid overfitting.
\end{itemize}
\end{propertybox}
